% Source : cours/cm7.qmd. Désinflation, pas nécessairement déflation.
\begin{board}{7}{Taux plus élevés : une transmission par la demande}
 {La BCE peut freiner la demande ; elle ne produit pas d'énergie.}
\stage{5.4}{1}{Taux directeurs en hausse}
\stage{16}{2}{Le crédit devient plus cher}
\stage{26.6}{3}{Moins de pression sur les prix}
\bank{5.4}{8.3}
\noteat{5.4}{7.2}{BCE}
\draw[flow] (8,8.5)--(8,12.4);
\captext{5.4}{La hausse se transmet\\aux conditions des banques.}
\house{13.6}{8.2}
\factory{18.2}{7.1}
\captext{16}{Achats et investissements\\peuvent être reportés.}
% Des flèches toutes positives : le niveau des prix ne baisse pas forcément.
\draw[flow] (23.5,7.5)--(23.5,12);
\draw[flow] (26.5,7.5)--(26.5,10.8);
\draw[flow] (29.5,7.5)--(29.5,9.4);
\draw[drawline] (22.7,7.2)--(30.2,7.2);
\noteat{26.6}{6.35}{Avec un délai}
\captext{26.6}{Inflation ralentie,\\activité freinée.}
\end{board}
